\documentclass[10pt,a4paper]{amsart}
\usepackage[margin=19mm]{geometry}
\usepackage{amsmath,amssymb,mathtools}
\usepackage[scr=rsfs,cal=euler]{mathalfa}
\usepackage{xfrac}
\usepackage[unicode,hidelinks,bookmarks=false]{hyperref}
\newcommand{\ie}{i.e.\ }
\newcommand{\cf}{cf.\ }
\newcommand{\resp}{respectively\ }
\newcommand{\Subsection}{\subsection}
\providecommand{\nobreakdash}{\nobreak\hbox{-}}
\setlength{\emergencystretch}{2em}
\multlinegap0pt
\usepackage{bm}
\usepackage{bbm}
\usepackage{dsfont}
\usepackage{xspace}
\usepackage{cancel}
\usepackage{mathtools}
\usepackage{amsthm}
\makeatletter
\@ifundefined{theorem}{\newtheorem{theorem}{Theorem}}{}
\@ifundefined{proposition}{\newtheorem{proposition}[theorem]{Proposition}}{}
\makeatother
\newcommand{\C}{\mathcal{C}}
\newcommand{\Hs}{H_{\mathrm{FW}}}
\newcommand{\R}{\mathbb{R}}
\newcommand{\T}{\mathbb{T}}
\newcommand{\dd}{\mathrm{d}}
\newcommand{\diag}{\operatorname{diag}}
\title[Rare Fluctuations from Normally Hyperbolic Invariant Manifol]{Rare Fluctuations from Normally Hyperbolic Invariant Manifolds}
\author{Stephen Wiggins}
\date{Source: arXiv:2608.19814v1 (CC BY 4.0)}
\begin{document}
\maketitle

\noindent\emph{Source and licence.} Rare Fluctuations from Normally Hyperbolic Invariant Manifolds. Version arXiv:2608.19814v1 (CC BY 4.0), distributed under the Creative Commons Attribution 4.0 International licence, as recorded in the arXiv licence metadata for this exact version and shown on the abstract page https://arxiv.org/abs/2608.19814v1. Full rights notice: licenses/arxiv-2608.19814-CC-BY-4.0.txt.\par\smallskip

\noindent\textit{Editorial scope.} Counted unit: the Proposition whose source label is \texttt{prop:counterexample} in the article (source0.tex lines 692--695, source label \texttt{prop:counterexample}), delivered together with the article's own complete original proof of it (source0.tex lines 697--765, one \texttt{proof} environment of 69 lines). No mathematics is written, repaired or re-derived: statement and proof are the article's own text line for line. Required context retained uncounted: none. Every definition, display and label of those lines is retained unchanged. Omitted from this package and not redistributed: 1--691, 696--696, 766--1846. Nothing inside a retained excerpt is omitted or altered: every retained body is delivered exactly as the article's lines are written, so no omission receipt of any kind accompanies this package. Citations: the retained excerpts carry no citation command at all, so no bibliography entry of the article is transcribed and no reference is redistributed. Reference adaptation: none was needed and none was made; every reference inside the delivered unit resolves inside it, so no unresolved reference or citation is produced. Printed slips kept literal: no printed word, symbol, hypothesis or number of the counted statement or of its proof was altered. Layout and class: the article's own class is \texttt{article} with packages including geometry, amsmath, amssymb, amsthm, mathtools, bm, graphicx, booktabs, array, microtype, enumitem, url, hyperref, cleveref; the wrapper is the lane's pinned \texttt{amsart} editorial wrapper at 10pt on A4, which restates the theorem-like environments the retained text needs and adds the article's own macro definitions that the retained text needs (\texttt{\textbackslash C}, \texttt{\textbackslash Hs}, \texttt{\textbackslash R}, \texttt{\textbackslash T}, \texttt{\textbackslash dd}, \texttt{\textbackslash diag}), with extra wrapper packages (bm, bbm, dsfont, xspace, cancel, mathtools). The delivered numbering is the unit's own. Assets: no retained span contains a figure environment, a graphics-inclusion command, a PDF-page inclusion, a table or a TikZ picture; no image file is copied into this package, the package declares no asset dependency and no image file is redistributed with it. Licence: version arXiv:2608.19814v1 of this article is distributed under the Creative Commons Attribution 4.0 International licence, as recorded in the arXiv licence metadata for this exact version and shown on the abstract page https://arxiv.org/abs/2608.19814v1. A full rights notice is delivered at licenses/arxiv-2608.19814-CC-BY-4.0.txt. Characters: every retained excerpt is pure ASCII, so the delivered engine, XeLaTeX with its default Latin Modern text font, reports no missing character.
\begin{proposition}[Normal hyperbolicity need not be inherited within the zero-energy hypersurface]
\label{prop:counterexample}
There exists a compact normally attracting deterministic invariant torus $N$ and a positive-definite diffusion tensor for which an exactly invariant symplectic center manifold has a compact regular zero-energy shell that is not normally hyperbolic as an invariant manifold of the zero-energy hypersurface $\Hs^{-1}(0)$.
\end{proposition}

\begin{proof}
Let $M=\T^2\times\R$ with coordinates $(\theta_1,\theta_2,z)$, and consider
\begin{equation}
\dot\theta_1=1,
\qquad
\dot\theta_2=0,
\qquad
\dot z=(-1+2\cos\theta_1)z.
\label{eq:counter-drift}
\end{equation}
The torus $N=\T^2\times\{0\}$ is invariant.  Along a deterministic orbit, $\theta_1(t)=\theta_1(0)+t$ and
\[
z(t)=z(0)\exp\!\left[-t+2\sin(\theta_1(0)+t)-2\sin\theta_1(0)\right].
\]
Since the sine terms remain bounded,
\[
|z(t)|\le e^4e^{-t}|z(0)|.
\]
Thus perturbations normal to the torus contract exponentially at asymptotic rate one, while the two tangent directions have zero exponential rate.  Hence $N$ is normally attracting.

Choose the constant positive-definite diffusion tensor $a=\diag(1,1,q)$, $q>0$.  With conjugate variables $(\eta_1,\eta_2,\zeta)$, the Freidlin--Wentzell Hamiltonian is
\begin{equation}
\Hs
=
\eta_1+(-1+2\cos\theta_1)z\zeta
+\frac12(\eta_1^2+\eta_2^2+q\zeta^2).
\label{eq:counter-H}
\end{equation}
The four-dimensional manifold $\C_0=\{z=\zeta=0\}$ is exactly invariant and symplectic.  On it the Hamiltonian reduces to
\[
h(\eta_1,\eta_2)=\eta_1+\frac12(\eta_1^2+\eta_2^2),
\]
so the zero-energy shell is
\begin{equation}
(\eta_1+1)^2+\eta_2^2=1,
\qquad
\Sigma_{0,0}\simeq\T^2\times S^1.
\label{eq:counter-shell}
\end{equation}
This shell is compact.  It is also regular because the gradient of $h$ with respect to $(\eta_1,\eta_2)$ is $(1+\eta_1,\eta_2)$, which cannot vanish on the circle \eqref{eq:counter-shell}.

Now select the point on the momentum circle with $\eta_1=-1$ and $\eta_2=1$.  The Hamiltonian equations restricted to $\C_0$ give
\[
\dot\theta_1=1+\eta_1=0,
\qquad
\dot\theta_2=\eta_2=1.
\]
Thus the Hamiltonian trajectory keeps $\theta_1$ constant.  Choose the constant value $\theta_1=\pi/3$.  At this value $-1+2\cos\theta_1=0$.  The transverse linearized equations along the selected Hamiltonian trajectory are therefore
\begin{equation}
\dot z=q\zeta,
\qquad
\dot\zeta=0,
\label{eq:nilpotent-transverse}
\end{equation}
or
\[
\frac{\dd}{\dd t}\begin{pmatrix}z\\ \zeta\end{pmatrix}
=
\begin{pmatrix}0&q\\0&0\end{pmatrix}
\begin{pmatrix}z\\ \zeta\end{pmatrix}.
\]
The transverse matrix is nilpotent: its square is zero.  Consequently
\[
e^{tA_\perp}=I+tA_\perp
=
\begin{pmatrix}1&qt\\0&1\end{pmatrix}.
\]
Transverse perturbations therefore grow at most linearly rather than exponentially, and both transverse Lyapunov exponents are zero.  There is no exponential transverse contraction or expansion along this trajectory.  Because the $(z,\zeta)$ directions are tangent to $\Hs^{-1}(0)$ along the shell and normal to the shell within that hypersurface, normal hyperbolicity of the zero-energy shell fails there.
\end{proof}

\end{document}
